\documentclass{article}%
\usepackage[T1]{fontenc}%
\usepackage[utf8]{inputenc}%
\usepackage{lmodern}%
\usepackage{textcomp}%
\usepackage{lastpage}%
\usepackage{geometry}%
\geometry{margin=2cm}%
\usepackage{expex}%
\usepackage[T1]{fontenc}%
\usepackage[utf8]{inputenc}%
\usepackage[spanish]{babel}%
%
\title{Análisis Lingüístico: Gramatica Normativa Kiche}%
\author{Generado por el TFM de Elías Carrasco}%
%
\begin{document}%
\normalsize%
\maketitle%
\section{Page 55}%
\label{sec:Page55}%
\subsection{c) Sintagma Estativo}%
\label{subsec:c)SintagmaEstativo}%

%
Este tipo de sintagma puede tener por núcleo un sustantivo, adjetivo, número, adjetivo %
posicional, adverbio o existencial. Dichas palabras pueden ser modificadas por %
elementos facultativos, como las partículas de negación, de interrogación y direccionales. %
Ejemplos: %
\subsection{Interrogación + adjetivo posicional + partícula}%
\label{subsec:Interrogacin+adjetivoposicional+partcula}%

%
\par\medskip%
\ex
\textit{¿Está sentado el presidente?} \\
\begingl
\gla ¿La t’uyul k’u ri k’amal b’e? //
\glb  //
\glc  //
\glft \textbf{Sintaxis:} [SET]\textsubscript{SET} //
\endgl
\xe
%
\medskip%
\subsection{c) Sintagma Adjetival}%
\label{subsec:c)SintagmaAdjetival}%

%
Este tipo de sintagma tiene como núcleo un adjetivo y puede llevar otro adjetivo más un %
intensificador. %
\par\medskip%
\ex
\textit{La naranja ya está agria.} \\
\begingl
\gla Ch'am chi ri alanxax. //
\glb  //
\glc  //
\glft \textbf{Sintaxis:} [SADJ]\textsubscript{SADJ} //
\endgl
\xe
%
\medskip%
\par\medskip%
\ex
\textit{Un palo bonito mojado está en el camino.} \\
\begingl
\gla Jun je’lalaj ch’aqal che' k'o pa b’e. //
\glb  //
\glc  //
\glft \textbf{Sintaxis:} [SADJ]\textsubscript{SADJ} //
\endgl
\xe
%
\medskip%
\subsection{Interrogación}%
\label{subsec:Interrogacin}%

%
\par\medskip%
\ex
\textit{¿Es gruesa la tabla?} \\
\begingl
\gla ¿La pim k’u le tz'alam? //
\glb  //
\glc  //
\glft \textbf{Sintaxis:} [SADJ]\textsubscript{SADJ} //
\endgl
\xe
%
\medskip%
\subsection{Afirmación}%
\label{subsec:Afirmacin}%

%
\par\medskip%
\ex
\textit{Sí, los tamales están cocidos} \\
\begingl
\gla Je', chaq' ri wa. //
\glb  //
\glc  //
\glft \textbf{Sintaxis:} [SADJ]\textsubscript{SADJ} //
\endgl
\xe
%
\medskip%
\subsection{e) Sintagma Adverbial}%
\label{subsec:e)SintagmaAdverbial}%

%
Este tipo de sintagma lleva un adverbio como núcleo, que puede ser modificado por uno %
o dos adverbios, partículas de interrogación y negación. %
\par\medskip%
\ex
\textit{Ayer tapizcamos.} \\
\begingl
\gla Iwir xujjach'anik. Ayer tapizcamos. //
\glb  //
\glc  //
\glft \textbf{Sintaxis:} [SAD]\textsubscript{SAD} //
\endgl
\xe
%
\medskip%
\par\medskip%
\ex
\textit{Mañana temprano, se van a pastorear.} \\
\begingl
\gla Nimaq'ab' chwe'q kixb'e pa yuq'. //
\glb  //
\glc  //
\glft \textbf{Sintaxis:} [SAD]\textsubscript{SAD} //
\endgl
\xe
%
\medskip%
\subsection{Adverbio + Partícula}%
\label{subsec:Adverbio+Partcula}%

%
\par\medskip%
\ex
\textit{¿Fue el año pasado que te casaste con la señorita Nikte'?} \\
\begingl
\gla ¿La junab'ir xatk'uli' ruk' ali Nikte'? //
\glb  //
\glc  //
\glft \textbf{Sintaxis:} [SAD]\textsubscript{SAD} //
\endgl
\xe
%
\medskip

%
\end{document}